\begin{helper}
Let \(D\) be a loopless \(T_s\)-free digraph on a vertex set \(W\).  Let
\(\ell,n\) be natural numbers with \(1\le\ell\), and let
\(\phi_c:[n]\to W\), for \(c<\ell-1\), be maps.  For every first color
\(c<\ell-1\) and every strictly increasing tuple
\[
  v:[s]\to[n],
\]
the coloring \(\noderef{RandomHomomorphismColoring}\) does not assign color
\(c\) to every edge \(\{v_i,v_j\}\) with \(i<j\).
\end{helper}
\begin{proof}
Suppose, for contradiction, that a first color \(c<\ell-1\) and a strictly
increasing tuple \(v:[s]\to[n]\) have every edge \(\{v_i,v_j\}\), \(i<j\),
colored \(c\).  Fix \(i<j\).  Since \(v\) is strictly increasing, the ordered
endpoints of the unordered edge \(\{v_i,v_j\}\) are \(v_i\le v_j\).  By the
definition of \(\noderef{RandomHomomorphismColoring}\), an edge can receive a
first color \(c\) only when the corresponding arc
\[
  D(\phi_c(v_i),\phi_c(v_j))
\]
is present.  Hence these arcs are present for all \(i<j\).

The map \(i\mapsto \phi_c(v_i)\) is injective.  Indeed, if
\(\phi_c(v_i)=\phi_c(v_j)\) with \(i\ne j\), then either \(i<j\) or \(j<i\).
In the first case the arc \(D(\phi_c(v_i),\phi_c(v_j))\) becomes a loop at
\(\phi_c(v_i)\), and in the second case the arc
\(D(\phi_c(v_j),\phi_c(v_i))\) becomes the same kind of loop.  Both cases
contradict the loopless hypothesis \(\noderef{DigraphLoopless}\).

Thus \(i\mapsto\phi_c(v_i)\) is an injective ordered \(s\)-tuple of vertices,
and the preceding paragraph gives all arcs from earlier indices to later
indices.  This is exactly a copy of the transitive tournament \(T_s\) in
\(D\), contradicting \(\noderef{TransitiveTournamentFree}\).  Therefore no
such monochromatic ordered \(s\)-clique in a first color exists.
\end{proof}
